\section{A finite connected-sum conflict lower bound}
\label{sec:conflict-composition}

The boundary counterexamples and the certified collar simplification leave
open whether a useful small resolving-set theorem holds on simplified knot
exteriors.  We next retain a diagram-authenticated connected-sum example
whose positive root-ray conflict hypergraph needs at least five tetrahedra
to hit.  The maintained optional inverse layering leaves this particular
source unchanged.  This is a finite obstruction to a universal bound of
four for these sources.  It is neither a growing family nor a lower bound
on recognition time, and it does not settle whether a polylogarithmic
bound holds after a stronger justified simplification.

\subsection{Exact diagram and source authentication}

Start with the closure of the three-strand braid
\[
 (\sigma_1)^3(\sigma_2)^3.
\]
The retained bare PD passes native diagram validation.  Its interlacement
factor certificate gives the crossing blocks $\{0,1,2\}$ and
$\{3,4,5\}$.  Each factor is identified with the same trefoil PD by an
explicit bijection of edge labels that preserves crossing order and all
four crossing slots.  No reflection or under/over exchange is used.
Thus the experiment starts from a visible diagrammatic connected sum of
two trefoils, rather than inferring that type from a polynomial invariant.

The experiment calls \path{recognize_by_crushing} directly on this entire
connected-sum diagram.  It does not describe the default high-level
\path{recognize} router: that interface already applies certified visible
interlacement factorization before its factor engines, as explained in
\path{research_updates.tex}.  The existing router can therefore avoid
constructing this combined cone.  Neither the source parameter bound nor
the failed whole-source candidates below are a lower bound on that router.

The native diagram exterior has $120$ tetrahedra.  The verified shelling
prefix leaves $114$, and $24$ independently verified vertex-reduction
crushes leave a finite, connected, orientable, one-vertex source $T$ with
one torus boundary and $23$ tetrahedra.  Its exact labelled-source digest
is
\begin{quote}\small\ttfamily
825247bfc387213d7cc965a33018a448043469eecf0a9edd3ce8d7a1b5fcb6f9.
\end{quote}
The retained \path{diagram-crushing-v1} certificate authenticates the
whole sequence from the supplied PD.  It has status \path{REDUCED}; no
recognition verdict is inferred from this partial chain.

A fresh control enables inverse layering and uses conflict-directed
transversal search, with zero tree-node, pivot and anchor allowances.
Vertex-reduction preprocessing still runs
before Euler search.  The resulting independently checked
\path{diagram-crushing-v2} chain contains exactly the same $24$ normal
crush packets, and its final labelled triangulation equals $T$ literally.
It performs zero inverse layers and no LP calls.  The boundary faces are
$(17,0)$ and $(20,0)$, in different tetrahedra.  Consequently there is no
candidate for the exposed two-boundary-face operation of
Section~\ref{sec:boundary-unlayer}.  This control proves survival through
the currently implemented optional preprocessing.  It gives no
invariance claim for $h$ under arbitrary triangulation moves or under
later positive normal crushing.

\subsection{Original STANDARD rank certificates for the root rays}

The evidence retains $35$ distinct primitive integral Q vectors with
canonical nonnegative formal standard lifts.  They all have positive
formal Euler characteristic and at least one incompatible tetrahedron.
No claim that these formal vectors are embedded surfaces is used.
Their positive coordinates were discovered by bounded native exact LP
queries; the fresh replay instead checks their original equations and
ranks without invoking an LP backend or Q-cycle elimination.

Let $M$ be the original standard matching matrix for $T$, and let $S$
be the positive Q support of a retained vector.  Keep \emph{all} $92$
triangle columns together with the $|S|$ Q columns.  Let $M_S$ be this
restricted matrix and $M_\triangle$ its triangle block.  The replay
checks
\[
 \operatorname{rank}(M_\triangle)=91,\qquad
 \operatorname{rank}(M_S)=90+|S|.
\]
In particular,
\[
 \dim\ker M_S=2,\qquad \dim\ker M_\triangle=1.
\]
Projection to the Q coordinates therefore has a one-dimensional image
on this supported matching space.  The triangle-only kernel is the
single vertex-link direction.  Conversely every vector in that image
has a standard matching lift; adding a sufficiently large multiple of
the finite vertex link makes its triangle coordinates nonnegative.
The supplied positive Q vector belongs to this image and uses every
column of $S$.  Its nonnegative supported cone is thus exactly one ray.
Since setting the other Q coordinates to zero defines a coordinate
face of the full cone, this ray is a \emph{root} extreme ray by
Lemma~\ref{lem:transversal-face-rays}.

The replay additionally checks every original standard matching equation,
nonnegative integer coordinates, Q gcd equal to one, and minimum triangle
coordinate zero.  Because the source has one boundary vertex, this last
condition is precisely canonical peeling of its vertex-link disc.
Formal Euler is evaluated directly as vertices minus arcs plus normal
pieces, using reconstructed global edge weights and face incidences.
It is not accepted from the discovery log or deduced from a knot-type
claim.  All $35$ vectors have $F>0$.

\subsection{Seven rays force five resolving tetrahedra}

The following seven retained rays already suffice.  Labels are
zero-based indices in the immutable $35$-ray list.  The table gives their
formal Euler value, positive Q support size, and complete conflict set.
\begin{center}
\begin{tabular}{rrrrl}
Ray & $F$ & $|S|$ & $\operatorname{rank}M_S$ & $D(r)$\\
\hline
$5$  & $1$ & $11$ & $101$ & $\{3,4\}$\\
$6$  & $1$ & $8$  & $98$  & $\{5\}$\\
$7$  & $1$ & $12$ & $102$ & $\{7,12,17,21\}$\\
$11$ & $1$ & $12$ & $102$ & $\{10,11,14\}$\\
$17$ & $2$ & $16$ & $106$ & $\{12,14\}$\\
$19$ & $1$ & $11$ & $101$ & $\{7,14,18,19\}$\\
$20$ & $2$ & $17$ & $107$ & $\{11,16,19\}$
\end{tabular}
\end{center}

\begin{proposition}[Finite authenticated lower bound]
\label{prop:connected-sum-h-five}
For the labelled source $T$ above, the uniform incompatible positive
root-ray conflict parameter of Section~\ref{sec:euler-transversal}
satisfies $h(T)\ge5$.  The same bound holds on the literally identical
endpoint of the retained optional inverse-layer preprocessing control.
\end{proposition}
\begin{proof}
Every listed ray is independently verified to be a positive incompatible
root extreme ray.  Hence a transversal for the complete root-ray
hypergraph must hit all seven listed conflict sets.

The first two sets require two separate choices, and their union is
disjoint from the last five sets.  Write those last five sets as
$A,B,C,D,E$ in table order.  Suppose two choices could hit all five.
Since $C=\{12,14\}$, one choice must be $12$ or $14$.
If it is $12$, that choice belongs to none of $B,D,E$, so the other
choice must lie in
\[
 B\cap D\cap E
 =\{10,11,14\}\cap\{7,14,18,19\}\cap\{11,16,19\}
 =\varnothing.
\]
If it is $14$, that choice belongs to neither $A$ nor $E$, so the
other must lie in
\[
 A\cap E=\{7,12,17,21\}\cap\{11,16,19\}=\varnothing.
\]
Both cases are impossible.  The last five sets therefore require at
least three further choices, proving $h(T)\ge2+3=5$.  The preprocessing
claim follows from verified equality of the labelled endpoints.
\end{proof}

For an independent finite enumeration check, every one of the
$\binom{23}{4}=8{,}855$ four-element subsets of all source tetrahedron
labels is compared with all retained conflict sets and with the
seven-set core.  Each misses at least one set.  Smaller subsets extend
to four-element subsets, so none can hit all the sets either.  The
retained set
\[
 B_0=\{3,5,7,11,14\}
\]
hits all $35$ observed conflict sets.  Thus the \emph{retained finite
hypergraph} has transversal number exactly five.  This is not an upper
bound of five for $h(T)$, since other positive root rays have not been
excluded.  A separate four-ray disjoint conflict packing is retained,
but yields only the weaker bound four.

\subsection{Discovery limits, replay and the open composition theorem}

The first full-source transversal query was inconclusive after its
$5{,}000$-pivot allowance.  Two deterministic $20$-query anchor batches
and five bounded explicit-candidate attempts supplied the retained
rays.  Each explicit candidate was refuted by a positive root ray
whose conflict set missed that candidate.  The last tested candidate,
$\{3,5,7,11,12\}$, was refuted by $\{14,18,19\}$.
No candidate attempt completed a negative cover.  In particular,
\emph{the canonical admissible cone of this source is not certified
negative}.  The knot type does not imply that property: a triangulation
of a nontrivial knot exterior may still contain positive canonical
normal spheres or boundary-parallel discs.  Such a compatible positive
witness could justify a later strict crushing step.  The finite result
therefore does not establish a difficult terminal source or a lower
bound on the work needed to recognize the connected sum.

The package \path{fast/conflict_composition_research} retains the exact
PD, factor maps, all normal coordinates, both complete source chains,
all exploratory configurations and outcomes, and the historical native
source hashes for each generation phase.  The original JSON bytes are
stored under deterministic gzip compression and authenticated by a
separate provenance manifest.  Fresh replay pins cover all current
native Python files and every immutable package input and source;
historical hashes are not substituted for those fresh pins.  The
solver-free audit and focused mutation controls are retained separately
in \path{synthesis/data/conflict-composition-*}.  No elapsed discovery
value is presented as a performance benchmark.

A future growing-family construction needs more than copying these
conflicts into several diagrams.  For a proposed gluing, let $V$ be the
original supported standard matching space and $L$ its triangle-only
kernel.  A sufficient rank-transport argument must control the induced
map on $V/L$, so that each copied Q support still has dimension one.
It must also verify the canonical Euler correction after the global
vertex links are peeled.  Triangle-only extension across an annular
port is constrained by the actual trace space of $L$; arbitrary
boundary arcs do not extend with zero Q coordinates.  A weaker useful
condition is \emph{conflict locality}: extensions may use Q coordinates
in other blocks, provided at most one type is used in each such
tetrahedron.  Disjoint local conflicts can then survive even when Q
supports overlap.  None of these transport conditions is inferred
from the connected-sum knot type.

The archived composition memo spells out these rank, trace and Euler
obligations.  It also distinguishes sources with an easy compatible
positive witness from sources with a completely certified negative
canonical cone.  A family that merely adds crushable sphere gadgets
could have growing source $h$ while still being simplified quickly.
No additive lower-bound family, componentwise upper-bound theorem,
simplification-survival theorem, or general quasi-polynomial recognition
bound follows from the present finite experiment.
